\begin{ThreePartTable}
\begin{longtable}[c]{|c|c|c|L{9cm}|}
    \caption{Instruction Field Names and Descriptions}
    \label{tab:pie-instruction-field}\\
    \hline
    \rowcolor{lightgray}
    \multicolumn{3}{|c|}{\textbf{Name}} & \textbf{Description}\\\hline
    \endfirsthead
    \multicolumn{4}{c}{\bfseries \tablename\ \thetable{} -- cont'd from previous page} \\\hline
    
    \rowcolor{lightgray}
    \multicolumn{3}{|c|}{\textbf{Name}} & \textbf{Description}\\\hline
    \endhead
    \multicolumn{4}{r}{\textbf{Cont'd on next page}} \\
    \endfoot
    \endlastfoot
    \multirow{6}{*}{a*} \label{var:pie_ar_reg_def} & \multirow{6}{*}{32-bit general-purpose registers} & as & In-out type (used as input/output operand). Stores address information for read/write operations, which is updated after such operations are completed.\\\cline{3-4}
            & & at & In-out type. Temporarily stores operation results to the \hyperref[instruction:EEFFTAMSS16STINCP]{EE.FFT.AMS.S16.ST.INCP} instruction, which will be part of the data to be written to memory.
            \\\cline{3-4}
            & & ad & In type (used as input operand). Stores data used to update address information. \\\cline{3-4}
            & & av & In type. Stores data to be written to memory. \\\cline{3-4}
            & & ax,ay & In type. Stores data involved in arithmetic operations, e.g., shifting amounts, multipliers and etc.
            \\\cline{3-4}
            & & au & Out type (used as output operand). Stores results of instruction operations.
            \\\hline
    \multirow{5}{*}{q*} & \multirow{5}{*}{128-bit general-purpose registers} & qs & In type. Stores 128-bit data used for concatenation operations.\\\cline{3-4}
            & & qa,qx,qy,qm & In type. Stores data used for vector operations.\\\cline{3-4}
            & & qz & Out type. Stores results of vector operations.\\\cline{3-4}
            & & qu & Out type. Stores data read from memory. \\\cline{3-4}
            & & qv & In type. Stores data to be written to memory. \\\hline
    \multicolumn{3}{|c|}{fu} & 32-bit general-purpose floating-point register. Stores floating-point data read from memory. \\\hline
    \multicolumn{3}{|c|}{fv} & 32-bit general-purpose floating-point register. Stores floating-point data to be written to memory. \\\hline
    \multicolumn{3}{|c|}{sel2} & 1-bit immediate value ranging from 0 to 1. Used to select signals. \\\hline
    \multicolumn{3}{|c|}{sel4,upd4} & 2-bit immediate value ranging from 0 to 3. Used to select signals. \\\hline
    \multicolumn{3}{|c|}{sel8} & 3-bit immediate value ranging from 0 to 7. Used to select signals. \\\hline
    \multicolumn{3}{|c|}{sel16} & 4-bit immediate value ranging from 0 to 15. Used to select signals. \\\hline
    \multicolumn{3}{|c|}{sar2} & 1-bit immediate value ranging from 0 to 1. Represents shifting numbers.  \\\hline
    \multicolumn{3}{|c|}{sar4} & 2-bit immediate value ranging from 0 to 3. Represents shifting numbers. \\\hline
    \multicolumn{3}{|c|}{sar16} & 4-bit immediate value ranging from 0 to 15. Represents shifting numbers. \\\hline
    \multicolumn{3}{|c|}{imm1} & 7-bit unsigned immediate value ranging from 0 to 127 with an interval of 1. This is used to show the size of the updated read/write operation address value. \\\hline
    \multicolumn{3}{|c|}{imm2} & 7-bit unsigned immediate value ranging from 0 to 254 with an interval of 2. This is used to show the size of the updated read/write operation address value. \\\hline
    \multicolumn{3}{|c|}{imm4} & 8-bit signed immediate value ranging from -256 to 252 with an interval of 8. This is used to show the size of the updated read/write operation address value. \\\hline
    \multicolumn{3}{|c|}{imm16} & 8-bit signed immediate value ranging from -2048 to 2032 with an interval of 16. This is used to show the size of the updated read/write operation address value. \\\hline
    \multicolumn{3}{|c|}{imm16f} & 4-bit signed immediate value ranging from -128 to 112 with an interval of 16. This is used to show the size of the updated read/write operation address value. \\\hline
\end{longtable}
\end{ThreePartTable}